\scanpage{161}\mspage{16}
l'établir consiste certainement dans une description soigneuse des catégories de coefficients avec lesquels on travaille, aussi bien sur $X$ et sur $k$, et des foncteurs $\deleted{\unsure{\Delta}}D$ et $\deleted{\unsure{D}}\Delta$ dessus. Comme la définition d'une flèche est immédiate, une fois le travail de montage de la machine fait, la validité de la formule doit résulter sans difficulté des dévissages habituels, permettant de nous ramener au cas standart $F=\mathbb{Z}/p\mathbb{Z}$, $\mu_p$, $\alpha_p$ sur $X$ lisse complète (\uline{NB} le cas des coeff. premiers à la car. est déjà connu.)

\noindent\tikz[baseline=-0.1ex,overlay]{\draw (-.43,0)--(-.17,0)--(-.17,.37)--(.53,.37);}Explicitant ce que donne la formule de \unsure{dualité} pour $R^1f_*(G_X)$, \unsure{où} $G$ est un groupe \inserted{fini} étale sur $k$ (le cas le plus important étant $G=(\mathbb{Z}/p^{\unsure{n|m}}\mathbb{Z})_k$), on retrouve la description à la Serre du \deleted{\unsure{``|le}}corps de classe ``géométrique'', en termes d'extensions par $G$ d'une jacobienne généralisée de $X$. \unsure{Ainsi}, la formule de dualité peut être comprise comme une version cohomologique considérablement enrichie du corps de \unsure{classes|classe} géométrique. Quand le corps de base $k$ est fini, pour retrouver le corps de \unsure{:} classe sous sa forme classique, on peut mettre ``l'astuce de Lang'' [sur la \unsure{relation|reltion} entre \deleted{\unsure{un|le|la}}\unsure{,} le $\pi_1$ ``arithmetique'' d'un groupe algébrique \inserted{comm.} lisse connexe \inserted{$J$} sur $k$ et son $\mathrm{H}^{0}(k,J)=J(k)$, -- le $\pi_1^{\mathrm{ar}}(J)$ \unsure{classifiant} les isogénies \inserted{au dessus de $J$} à noyau un \unsure{groupe} \unsure{con}
% archive sheet number: 161
